\documentclass[a4paper,oneside,onecolumn,article,draft]{memoir}
\counterwithout{section}{chapter}

\DoubleSpacing

\usepackage[T1]{fontenc}
\usepackage{palatino}
\usepackage[final]{graphicx}
\graphicspath{{./figs/}}

%% remove colorlinks when preparing for print
\usepackage[final,hyperindex,hyperfootnotes,bookmarksnumbered,colorlinks]{hyperref}

\usepackage{eqparbox}
%% keep size of first column of sequence align table on memory to calculate
%% the size limit for the second column
\newsavebox{\firstentrybox}
\newcolumntype{F}{%
  >{\begin{lrbox}{\firstentrybox}}%
    l%
  <{\end{lrbox}%
  \eqmakebox[firstentry][l]{\unhcopy\firstentrybox}}}

\newsubfloat{figure}

\usepackage{color}    % to color background of generated values values
\usepackage{seqsplit}
\usepackage{url}
\usepackage[textsize=footnotesize]{todonotes}
%% new command for box about missing references
\newcommand{\addref}[1][]{\todo[color=red!40,size=\tiny]{Add reference: #1}}

\usepackage[sectionbib]{natbib}

\usepackage{fp}
\usepackage{capt-of}

\newfloat{supplement}{sup}{S}
\setfloatlocations{supplement}{h}
\renewcommand{\thesupplement}{S\arabic{supplement}}

%% to keep a pair of related table and figures, we set a new float and use it
%%conjunction with the captionof command from capt-of
\newfloat{TableAndFigure}{X}{capname}

\input{results/variables-align_results}
\input{results/variables-cluster_stats}
\input{results/variables-protein_stats}
\input{results/variables-utr}
%% NOTE: we are not counting H1 for number of histones
%% NOTE: counted from Table 1

% 13 coding H2A, 9 coding H3, 6 coding H1
% 13 coding H2B, 10 coding H4
% 1 pseudo H2A (Hist1h2aj with the _ -- see caption of Fig 1)
\newcommand{\MarzluffCodingGenesInHistOne}{45}
\newcommand{\MarzluffPseudoGenesInHistOne}{1}
\newcommand{\MarzluffTotalGenesInHistOne}{46}

% 4 coding H2A, 0 pseudo H3, 3 coding H3
% 0 pseudo H2B, 2 coding H2B, 1 coding H4
\newcommand{\MarzluffCodingGenesInHistTwo}{10}
\newcommand{\MarzluffPseudoGenesInHistTwo}{0}
\newcommand{\MarzluffTotalGenesInHistTwo}{10}

% 1 coding H2A, 0 coding H3
% 0 pseudo H2B, 2 coding H2B
\newcommand{\MarzluffCodingGenesInHistThree}{3}
\newcommand{\MarzluffPseudoGenesInHistThree}{0}
\newcommand{\MarzluffTotalGenesInHistThree}{3}

% 1 coding H4
\newcommand{\MarzluffCodingGenesInHistFour}{1}
\newcommand{\MarzluffPseudoGenesInHistFour}{0}
\newcommand{\MarzluffTotalGenesInHistFour}{1}

\newcommand{\MarzluffHistOneSpan}{2.25\,Mbp} % from Fig.1 (B)
\newcommand{\MarzluffHistTwoSpan}{60\,kbp} % from Fig.2 (A)
\newcommand{\MarzluffHistThreeSpan}{8\,kbp} % from Fig.2 (B)
\newcommand{\MarzluffHistFourSpan}{}


%% We define a style here which will be used around all the values found
%% by the scripts. The idea is to place them in a grey box so that people
%% reading the paper can really get the idea that the values in the text
%% are automatically generated. The problem is that the variables are not
%% simple numbers or text, and we want to make operations with them here
%% in LaTeX. So we get an handy trick, and redefine FPeval.
%% See http://tex.stackexchange.com/questions/159155 for details
\newcommand{\ScriptValue}[1]{\setlength{\fboxsep}{2pt}\colorbox[gray]{0.8}{#1}}
\let\OldFPeval\FPeval
\def\FPeval{\renewcommand{\ScriptValue}[1]{##1}\OldFPeval}

\input{results/variables-reference_comparison}
% 12 coding H2A, 10 coding H3, 6 coding H1
% 15 coding H2B, 12 coding H4
\newcommand{\MarzluffCodingGenesInHistOne}{50}
\newcommand{\MarzluffPseudoGenesInHistOne}{9}
\newcommand{\MarzluffTotalGenesInHistOne}{59}

% 3 coding H2A, 2 pseudo H3, 1 coding H3
% 4 pseudo H2B, 1 coding H2B, 1 coding H4
\newcommand{\MarzluffCodingGenesInHistTwo}{11}
\newcommand{\MarzluffPseudoGenesInHistTwo}{7}
\newcommand{\MarzluffTotalGenesInHistTwo}{XXXXX}

% 1 coding H2A, 1 coding H3
% 1 pseudo H2B, 1 coding H2B
\newcommand{\MarzluffCodingGenesInHistThree}{3}
\newcommand{\MarzluffPseudoGenesInHistThree}{XXXX}
\newcommand{\MarzluffTotalGenesInHistThree}{XXXXX}

% 1 coding H2A
\newcommand{\MarzluffCodingGenesInHistFour}{1}
\newcommand{\MarzluffPseudoGenesInHistFour}{0}
\newcommand{\MarzluffTotalGenesInHistFour}{0}

\newcommand{\MarzluffHistOneSpan}{2.3\,Mbp}
\newcommand{\MarzluffHistTwoSpan}{105\,kbp}
\newcommand{\MarzluffHistThreeSpan}{XXXXbp}
\newcommand{\MarzluffHistFourSpan}{XXXXbp}


%%
%% Some computations fit better here in LaTeX than the scripts.
%%

\FPeval{\result}{\SequenceChangeSinceReference{}
                + \PseudoSinceReference{}
                + \CodingSinceReference{}
                + \AddedSinceReference{}
                + \RemovedSinceReference{}}
\FPround{\TotalChangesSinceReference}{\result}{0} % of course is an integer

\FPeval{\result}{\SequenceChangeSinceReference{}
                +\PseudoSinceReference{}
                +\CodingSinceReference{}}
\FPround{\SequencesUpdatedSinceReference}{\result}{0} % of course is an integer


\author{David Miguel Susano Pinto \and Andrew Flaus}
\title{The Human Canonical Core Histone Catalogue}

\begin{document}

  \maketitle

  \begin{abstract}
    Core histone proteins H2A, H2B, H3, and H4 are encoded
    by a large family of genes distributed across the human genome.
    Canonical core histones contribute the majority of proteins to bulk chromatin packaging,
    and are encoded in \NumberOfClusters{} clusters
    by \TotalCoreCodingGenes{} coding genes comprising
    \HTwoACodingGenes{} for H2A,
    \HTwoBCodingGenes{} for H2B,
    \HThreeCodingGenes{} for H3,
    and \HFourCodingGenes{} for H4,
    along with at least \TotalCorePseudoGenes{} total pseudogenes.
    The canonical core histone genes display coding variation that gives rise to
    \HTwoAUniqueProteins{} H2A, \HTwoBUniqueProteins{} H2B,
    \HThreeUniqueProteins{} H3, and \HFourUniqueProteins{} H4 unique protein isoforms.
    Although histone proteins are highly conserved overall,
    these isoforms represent a surprising and seldom recognised variation
    with amino acid identity as low as
    \FPround{\result}{\HallMinPID}{0} \SI{\result}{\percent}
    between canonical histone proteins of the same type.
    The gene sequence and protein isoform diversity
    also exceeds commonly used subtype designations such as H2A.1 and H3.1,
    and exists in parallel with the well-known specialisation of variant histone proteins.
    RNA sequencing of histone transcripts shows evidence for
    differential expression of histone genes
    but the functional significance of this variation has not yet been investigated.
    To assist understanding of the implications of histone gene and protein diversity
    we have curated and catalogued the entire human canonical core
    histone gene and protein complement.
    In order to organise this information in a
    robust, accessible, and accurate form,
    we have applied software build automation tools to
    dynamically generate the canonical core histone repertoire
    based on current genome annotations
    and then to organise the information into a manuscript format.
    Automatically generated values are shown with a light grey background.
    Alongside recognition of the encoded protein diversity,
    this has led to multiple corrections to human histone annotations,
    reflecting the flux of the human genome as it is updated and
    enriched in reference databases.
    This approach is inspired by the aims of reproducible research
    and can be readily adapted to other gene families.
  \end{abstract}

  \section{Introduction}

  Histones are among the most abundant proteins in eukaryotic cells
  and contribute up to half the mass of chromatin.
  The core histone types H2A, H2B, H3, and H4
  define the structure and accessibility of the nucleosome
  as the fundamental repeating unit of genome organisation
  around which the DNA is wrapped \citep{Luger1997structure}.
  In addition, the many chemically reactive sidechains of histones
  are post-translationally modified
  as a nexus for signalling and heritable epigenetics.

  Histones are delineated as either canonical or variant based on
  their gene location, expression characteristics
  and functional roles (\tref{tab:typical-histone-differences}).
  Canonical histones contribute the majority of proteins to
  the bulk structure and generic function of chromatin,
  and are encoded by \TotalGenes{} genes in \NumberOfClusters{}
  clusters in the human genome,
  of which \TotalCodingGenes{} are coding genes and \TotalPseudoGenes{}
  are pseudogenes (\tref{tab:histone-gene-count}).

  \begin{table}
    \caption{Properties distinguishing canonical and variant histone proteins.}
    \label{tab:typical-histone-differences}
    \centering
    \begin{tabular}{l l l}
      \toprule
      \null                     & Canonical             & Variants \\
      \midrule
      Expression timing         & Replication dependent & Replication independent \\
      Sequence identity         & High                  & Low \\
      Functional relationships  & Isoforms              & Specialised functions \\
      Transcript stabilisation  & Stem-loop             & poly(A) tail \\
      Gene distribution         & Clusters              & Scattered \\
      \bottomrule
    \end{tabular}
  \end{table}

  \begin{table}
    \caption{Count of human canonical histone coding and pseudogenes by
             histone cluster and type. $\psi$ indicates pseudo genes.}
    \label{tab:histone-gene-count}
    \centering
    \begin{tabular}{l *{5}{r!{+}r<{$\psi$}}}
      \toprule
      \null   & \multicolumn{2}{c}{H2A}  & \multicolumn{2}{c}{H2B}
              & \multicolumn{2}{c}{H3}   & \multicolumn{2}{c}{H4}
              & \multicolumn{2}{c}{Total} \\
      \midrule
      HIST1   & \HTwoACodingInHISTOne{}     & \HTwoAPseudoInHISTOne{}
              & \HTwoBCodingInHISTOne{}     & \HTwoBPseudoInHISTOne{}
              & \HThreeCodingInHISTOne{}    & \HThreePseudoInHISTOne{}
              & \HFourCodingInHISTOne{}     & \HFourPseudoInHISTOne{}
              & \CodingGenesInHISTOne{}     & \PseudoGenesInHISTOne{} \\
      HIST2   & \HTwoACodingInHISTTwo{}     & \HTwoAPseudoInHISTTwo{}
              & \HTwoBCodingInHISTTwo{}     & \HTwoBPseudoInHISTTwo{}
              & \HThreeCodingInHISTTwo{}    & \HThreePseudoInHISTTwo{}
              & \HFourCodingInHISTTwo{}     & \HFourPseudoInHISTTwo{}
              & \CodingGenesInHISTTwo{}     & \PseudoGenesInHISTTwo{} \\
      HIST3   & \HTwoACodingInHISTThree{}   & \HTwoAPseudoInHISTThree{}
              & \HTwoBCodingInHISTThree{}   & \HTwoBPseudoInHISTThree{}
              & \HThreeCodingInHISTThree{}  & \HThreePseudoInHISTThree{}
              & \HFourCodingInHISTThree{}   & \HFourPseudoInHISTThree{}
              & \CodingGenesInHISTThree{}   & \PseudoGenesInHISTThree{} \\
      HIST4   & \HTwoACodingInHISTFour{}    & \HTwoAPseudoInHISTFour{}
              & \HTwoBCodingInHISTFour{}    & \HTwoBPseudoInHISTFour{}
              & \HThreeCodingInHISTFour{}   & \HThreePseudoInHISTFour{}
              & \HFourCodingInHISTFour{}    & \HFourPseudoInHISTFour{}
              & \CodingGenesInHISTFour{}  & \PseudoGenesInHISTFour{} \\
      \addlinespace
      Total   & \HTwoACodingGenes{}       & \HTwoAPseudoGenes{}
              & \HTwoBCodingGenes{}       & \HTwoBPseudoGenes{}
              & \HThreeCodingGenes{}      & \HThreePseudoGenes{}
              & \HFourCodingGenes{}       & \HFourPseudoGenes{}
              & \TotalCodingGenes{}       & \TotalPseudoGenes{} \\
      \bottomrule
    \end{tabular}
  \end{table}

  Relationships within the histone family have been described using a variety of terminologies
  reflecting biochemical, functional and genomic perspectives that are briefly described below
  and summarised in \tref{tab:histone-divisions}.

  %% The following is a rather long table, longer than 1 page so it can't
  %% go in a float.  There are several packages such as longtable or xtab
  %% but those must be really used with tables.  However what we really
  %% have is a list of descriptions.  If we converted this into a table,
  %% like we had before, it would be a huge waste of space.  So we use
  %% afterpage and captionof to mimic a floating table.  And we use the
  %% shaded environment just because it looks good (although it does
  %% introduce an annoyingly long vertical space before the caption).
  \afterpage{
    \captionof{table}{Terminology describing histone variation}
    \label{tab:histone-divisions}
    \definecolor{shadecolor}{gray}{0.9}
    \begin{shaded}
      \begin{description}
        \item[Allelic variants] \hfill \newline
        Copies of a canonical histone type genes. 
        Not strictly alleles. Not histone variants.

        \item[Canonical histones]\hfill \newline
        Core histones with properties described in \tref{tab:typical-histone-differences}.
        Contribute the majority of histones in chromatin.
        Encompasses multiple protein isoforms.
        Complement of variant histones.

        \item[Core histones]\hfill \newline
        Histones that form part of the nucleosome core particle wrapping 147 bp of DNA,
        comprising histone types H2A, H2B, H3, and H4.
        Encompasses both canonical and variant histones, 
        and complements linker histones.

        \item[Heteromorphous variants] \hfill \newline
        Histone variants with distinct function and localisation
        that are readily separated by gel electrophoresis.

        \item[Homomorphous variants] \hfill \newline
        Canonical histone subtypes
        requiring high resolution separation methods such as TAU PAGE.

        \item[Families] \hfill \newline
        Synonym for ``histone types''.

        \item[Isoforms] \hfill \newline
        Proteins with high sequence identity and largely equivalent function.
        Functional equivalence has not been demonstrated for canonical histone isoforms.

        \item[Linker histones] \hfill \newline
        Histones binding to linker DNA adjacent to the nucleosome core particle.
        The two linker histone types are H1 and H5. Complement of core histones.

        \item[Non-allelic variants] \hfill \newline
        Synonym for ``variant histones''. See also ``allelic variants''.

        \item[Replacement histones] \hfill \newline
        Synonym for ``variant histones''.
        Named because they can replace canonical core histones assembled in S phase.

        \item[Replication-dependent histones] \hfill \newline
        Synonym for ``canonical histones''.
        Named because expression occurs primarily in S phase.

        \item[Replication-independent histones] \hfill \newline
        Synonym for ``variant histones''.
        Named because expression is not predominantly in S phase.

        \item[Subtypes] \hfill \newline
        Canonical histone type isoforms separable by TAU PAGE
        (e.g. H2A.1, H2A.2).  Despite their naming, there is no evidence of
        functional differences between them.

        \item[Types] \hfill \newline
        Histone proteins sharing sequence homology
        that participate in specific combinations to define the repeating nucleosome structure.
        The 5 histone types are H1, H2A, H2B, H3 and H4.

        \item[Variant histones] \hfill \newline
        Core histones with properties described in \tref{tab:typical-histone-differences}.
        Contribute a minor proportion of histones in chromatin and perform specialised functions.
        Complement of canonical histones.

        \item[Wild type histones] \hfill \newline
        Synonym of ``canonical histones''.
      \end{description}
    \end{shaded}
  }

  \subsection{Biochemical perspective}

    Abundant histone proteins are readily isolated using their
    highly basic chemical character.
    Successive improvements in fractionation ultimately revealed 5 main histone types
    with nomenclature H1, H2A, H2B, H3, and H4 \citep{nomenclature}.
    An additional H1--related histone H5 is recognised in avian erythrocytes \citep{HFive-review}.

    The demonstration of the nucleosome as the fundamental
    repeating unit of chromatin \citep{Kornberg1974}
    revealed that H2A, H2B, H3, and H4 associate as an octamer of two copies each within the
    nucleosome core particle. These four histones are referred to as core histones.
    In contrast, H1 associates with the linker DNA between nucleosome core particles
    and is referred to as a linker histone.
    The 5 human somatic H1 isoforms and 6 tissue-specific
    variants are described elsewhere \citep{HarshmanFreitas2013}.

    Arginine and lysine content was used as an early distinction between the histones.
    The H1 linker histone has a low arginine/lysine ratio
    of \LinkerArgLysRatio{} and became known as lysine-rich
    whereas the 4 core histones are arginine-rich
    with high arginine/lysine ratios of \HTwoAArgLysRatio{} in H2A, \HTwoBArgLysRatio{} in H2B,
    \HThreeArgLysRatio{} in H3, and \HFourArgLysRatio{} in H4 type isoforms.
    Nevertheless, the core histones contains many lysines particularly in their N-terminal tails.

    Separating histones by polyacrylamide gel electrophoresis (PAGE)
    using the strongly anionic detergent sodium dodecyl sulphate and neutral buffers (SDS PAGE)
    gives single bands for each histone type \citep{ShechterHake2007}.
    However, PAGE with non-ionic detergent Triton X--100 and uncharged urea as denaturants
    in acid buffers (TAU or AUT PAGE) allows the separation
    of histone types into multiple bands
    due to post-translational modifications and differences at specific amino acids
    in the polypeptide \citep{Zweidler1977}.
    These TAU PAGE separations gave rise to subtype designations
    H2A.1, H2A.2, H3.1, H3.2, and H3.3.

  \subsection{Functional perspective}

    Canonical core histone expression
    is significantly elevated during S~phase to provide chromatin packaging
    for DNA duplicated during replication.
    This led to their description as ``replication dependent'',
    although a supply of canonical histones is inevitably required
    to partner variants throughout the cell cycle.
    Metazoan canonical histone genes are distinctive
    because they lack introns and give rise to non-polyadenylated protein coding transcripts.
    Turnover is independently regulated via a highly
    conserved 3' stem-loop (\tref{tab:typical-histone-differences}).

    In contrast, variant histones such as H2A.Z, TH2B, and H3.3 have
    less sequence identity and lower abundance.
    They play functionally specific roles and are mostly expressed outside S~phase,
    so are described as ``replication independent''.
    Since histone variants are interpreted as taking the place
    of equivalent canonical histone types,
    they are also referred to as ``replacement'' histones.

  \subsection{Genomic perspective}

    Canonical histone genes are found in \NumberOfClusters{} clusters,
    possibly for regulatory efficiency.
    The multiple gene copies in these clustered arrays are
    sometimes confusingly referred to as ``allelic''
    and the resulting combined protein isoforms are often considered to be ``wild type''
    although both genes and protein products display
    variation in sequence and relative abundance,
    and their functional equivalence has not been tested.
    In contrast, almost all variant histones are encoded by single genes dispersed in the genome
    with typical properties including introns, alternative splicing,
    and polyadenylated transcripts (\tref{tab:typical-histone-differences}).

  \subsection{Canonical histone diversity}

    Despite the importance of histones for chromatin organisation and extensive interest
    in their role in epigenetics and regulation, the systematic
    curation and classification of human histone
    gene and protein sequences has not been revisited
    since the landmark survey of \citet{Marzluff02}.
    The canonical histone catalogue has accumulated
    \FPeval{result}{round(\SequenceChangeSinceReference{}
                          + \PseudoSinceReference{}
                          + \CodingSinceReference{}
                          + \AddedSinceReference{}
                          + \RemovedSinceReference{}, 0)}
    \result{}~differences (\tref{tab:difference-from-Marzluff02})
    from that survey as rich annotations continue to propagate
    into reference sequence databases.

    In this manuscript we provide a comprehensive catalogue
    of canonical histone genes, proteins, and pseudogenes.
    Since curation and annotation are dynamic and evolving,
    we have implemented the manuscript so that it can
    self-update from the most current data in the NCBI RefSeq database
    in order to maintain its value as a reference source in an accessible format.
    All automatically generated values here are displayed with a light grey background
    and were calculated using NCBI RefSeq data on \SequencesDate{}.
    \todo{Can't we generate linux timestamp and apply a formatting string in latex?}.
    This manuscript generation process has remained stable in our laboratory for several years
    and represents an example of ``reproducible research'' \citep{Claerbout2000}
    that provides a novel model for curating relatively stable gene families.

  \section{Histone genes}

  \subsection{Histone gene clusters}
	%% For simplicity, decided to list HIST4 explicitly for rather than allow it to come in by a variable
    The human canonical histone genes are located in clusters HIST1 to HIST4, 
	named in order of decreasing histone gene count (\tref{tab:histone-gene-count}).
	The major histone gene cluster HIST1 at chromosome~\HISTOneLocus{} and 
	contains \todo{percent genes in HIST1}\% of all histone genes 
	comprising \CodingGenesInHISTOne{}~functional core histone genes plus all canonical H1 histones. 
	HIST2 at chromosome~\HISTTwoLocus{} contains \CodingGenesInHISTTwo{}~coding genes. 
	HIST3 at \HISTThreeLocus{} and HIST4 at \HISTTFourLocus{} 
	contain only \CodingGenesInHISTThree{} and \CodingGenesInHISTFour{}~coding genes respectively. 
	Histone variants are dispersed across the genome outside the clusters, 
	so for example the three H3.3 encoding genes are located on chromosomes 1, 12 and 17 
	each far removed from other histone genes.

	HIST1 and HIST2 are both contiguous high density arrays of histone genes. 
	HIST1 spans \HISTOneSpan{} \todo{need space between number and unit}
	and is the second most gene dense region in the human genome at megabase scale \citep{MHC-III-analysis}.
	If \textit{HIST1H2AA} and \textit{HIST1H2BA} are excluded (see below), 
	the hemochromatosis protein encoding gene \textit{HFE} is the sole non-histone gene within HIST1. 
	There are no non-histone proteins within the HIST2 or HIST3 clusters. 

	The implications of this tight and exclusive clustering is not well understood. 
	The dispersive organisation has been suggested to imply that expression is regulated by local elements \citep{close-regulators}, 
	and it is also possible that clustering influences histone gene sequence conservation. 
	HIST1 is located towards the distal end of the major histocompatibility complex (MHC) in the extended class~I region \citep{MHC-I-transcript, MHC-complete-sequencing-1999} 
	and it has been suggested that proximity to HIST1 suppresses recombination in the MHC Class~I region \citep{MHC-repressed-by-HIST}. 
	It has also been reported that HIST2 is prone to deletions and frequent rearrangements \citep{HISTTwo-prone-deletion-discovery, HISTTwo-prone-deletion-focus}. 

	%% Removed birth and death because it is not directly relevant to this data
	%% Removed Alu stuff because it has no biological implications
	%% Definitely need to keep these observation, maybe put in thesis discussion?

  \subsection{Histone gene nomenclature}
    Canonical histone gene names adhere to a Human Genome Organisation (HUGO) Gene Nomenclature Committee (HGNC) 
	endorsed system derived from the cluster number and position relative to other histones \citep{Marzluff02}.
	%% Should we try putting a url reference into the manuscript?
	%% http://www.genenames.org/genefamilies/histones

	The human canonical histone gene symbols are divided into 3 parts: 
	HIST cluster, histone type, and identifier letter
	for the order relative to other histone genes of the same type in the cluster (\fref{fig:nomenclature}A).
	For example, \textit{HIST1H2BD} is the fourth H2B gene in the HIST1 cluster.
	%% Changed example to match figure!
	The direction of the cluster order varies, so HIST1 is ordered from the centromeric side 
	whereas HIST2 and HIST3 are ordered from the telomeric side.
	%% Think this is correct from viewing the genome
	
    Two exceptions were originally applied to this simple gene naming rule \citep{Marzluff02}. 
	Firstly, the position identifier is omitted if there are no other histones of the same type in the cluster, 
	so \textit{HIST3H2A} is the sole H2A gene in HIST3. 
	Secondly, the human and mouse histone clusters are largely syntenic 
	so positional identifiers for missing orthologues were skipped to maintain the equivalence of gene symbols.
	This means there is no human \textit{HIST1H2AF} although \textit{Hist1h2af} exists in mouse.

	Several new histone genes have been uncovered since the original naming (\tref{tab:difference-from-Marzluff02}) 
	and this required additional nomenclature exceptions.
	For example, new H2A encoding genes identified in the human and mouse HIST2 cluster after \textit{HIST2H2AA} 
	obliged renaming of \textit{HIST2H2AA} to \textit{HIST2H2AA3} and the addition of the new gene as \textit{HIST2H2AA4}. 
	There are no human orthologues of mouse \textit{Hist2h2aa1} and \textit{Hist2h2aa2}.
	
	Furthermore, no distinction was originally made between pseudogenes and functional coding genes, 
	so \textit{HIST3H2AA} is a pseudogene whereas its neighbour \textit{HIST3H2AB} is a coding gene.
	Newly uncovered pseudogenes are now suffixed with PS and a number in order of discovery 
	instead of a positional identifier, such as \textit{HIST1H2APS5} 
	However, the ``original'' pseudogenes retain their symbols without PS. 
	This means that the absence of a PS suffix does not indicate a functional gene, 
	and that there is no positional information in the gene symbols of new pseudogenes.

	The flux in genome sequencing and increasing interest in histone genes 
	leaves open the possibility of further change. 
	For example, \textit{HIST1H2BA} is now recognised as testes-specific TH2 \addref{Horokoshi ACTA Cryst 2014?} 
	and has the most divergent sequence of current canonical H2B proteins 
	and is ~300 kbp upstream of the rest of the HIST1 cluster 
	separated by a number of unrelated non-histone genes. 
	\textit{HIST1H2BA} and the adjacent divergent \textit{HIST1H2AA} therefore share some features with variant histones.
	
	All recognised histone variant gene familes comprise only one or a few copies 
	dispersed across the genome outside clusters.
	The histone variant gene symbols have a separate 3 part nomenclature 
	consisting of labels for histone type, F for family, 
	and an identifier letter based on discovery order (\fref{fig:nomenclature}B).
	Where two or more isoforms exist the copies gain a numerical suffix, 
	although this is sometimes omitted for the first copy. 
	For example, \textit{H2AY} and \textit{H2AY2} encode the human macroH2As, 
	with \textit{H2AY} giving rise to alternative transcripts encoding macroH2A1.1 and macroH2A1.2 protein isoforms 
	whereas \textit{H2AY2} encodes macroH2A2.
	The variant histone gene symbols (\tref{tab:histone-variant-families}) also contain exceptions. 
	For example, the three gene copies encoding the H3.3 variant 
	take on distinguishing letters \textit{H3F3A}, \textit{H3F3B} and \textit{H3F3C} instead of numbers 
	and the H3 variant gene CENPA was named as the first identified centromeric protein \cite{CENPA-first-report} 
	before it was realised to be a histone protein \citep{CENPA-copurifies-histones, CENPA-sequence-analysis}. 
	Histone variant pseudogenes are suffixed with P to yield symbols such as \textit{H2AFXP} and \textit{H2AFZP5}.
	The increasing interest in histone variant function \addref{Maze Allis NRG 2014}
	but imprecise usage and conflicts between species, 
	has led to guidelines for improved consistency \addref{Talbert Epigen Chrom 2012}.

    \begin{figure*}
      \centering
      \includegraphics[width=\textwidth]{nomenclature-schematic.pdf}
      \caption{Histone gene nomenclature. 
	  a) Canonical histone gene names encode their relative genomic order by cluster. 
	  b) Variant histone symbols are identified with F then a family identifier letter
	  and numbered by discovery, except 1 is omitted for the founder member. 
	  H3.3 genes conversely use numbers and letters for family and member.
	  c) Canonical seudogenes identified since 2002 are named by cluster with PS identifier numbered by discovery. 
	  Previously canonical pseudogenes were named like functional histone genes.}
      \label{fig:nomenclature}
	  \todo{Parts A and B need labels. Name text need to be italics!}
	  \todo{Change part B to use H2AFY2 as a more complete example of variant nomenclature}
	  \todo{Include part C using HIST1H2APS2 as example of PS pseudogene nomenclature}
    \end{figure*}

%% Let's leave the promoters and introns because we don't have anything to say

  \begin{table*}
    \caption{HGNC recognised histone variant family stem names, commonly used protein names 
	and protein names for improved consistency proposed by Talbert et al (2012). 
	HIST1H2BA is nominally not a variant but is included for completeness.}
    \label{tab:histone-variant-families}
    \centering
    \begin{tabular}{l l l}
      \toprule
      \Family & Protein name & Talbert et al & Notes \\
      \midrule
	  H2AFB & H2A.Bbd & H2A.B & equivalent to H2A.L \\
	  H2AFJ & H2A.J & ? & ? \\
	  H2AFV & H2A.Z-2 & H2A.Z.1 & - \\
	  H2AFX & H2AX & H2A.X & DNA damage related \\
	  H2AFY & macroH2A/mH2A & macroH2A & - \\
	  H2AFZ & H2A.Z & H2A.Z.1 & - \\
	  H2BFM & ? & ? & - \\
	  H2BFS & - & - & pseudogene \\
	  H2BFWT & ? & H2B.W & testes specific \\
	  H2BFX & - & - & pseudogene \\
	  HIST1H2BA & TH2B/TSH2B & H2B.1 & testes specific \\
	  H3F3 & H3.3 & H3.3 & - \\
	  CENPA & CENP-A & ? & centromere-specific \\
	\bottomrule
    \end{tabular}
  \end{table*}

%% This is arguably the most important table. We should cite it as early as possible
    \begin{table*}
      \caption{
	  Count of human canonical histone coding and pseudogenes by histone cluster and type.
	  The right-hand value in each cell marked with $\psi$ is the number of pseudo genes.}
      \label{tab:histone-gene-count}
      \centering
      \begin{tabular}{l r@{ + }l<{$\psi$} r@{ + }l<{$\psi$} r@{ + }l<{$\psi$} r@{ + }l<{$\psi$} r@{ + }l<{$\psi$}}
        \toprule
        \null   & \multicolumn{2}{c}{H2A}  & \multicolumn{2}{c}{H2B}
                & \multicolumn{2}{c}{H3}   & \multicolumn{2}{c}{H4}
                & \multicolumn{2}{c}{Total} \\
        \midrule
        HIST1   & \HTwoACodingInHISTOne{}     & \HTwoAPseudoInHISTOne{}
                & \HTwoBCodingInHISTOne{}     & \HTwoBPseudoInHISTOne{}
                & \HThreeCodingInHISTOne{}    & \HThreePseudoInHISTOne{}
                & \HFourCodingInHISTOne{}     & \HFourPseudoInHISTOne{}
                & \CodingGenesInHISTOne{}     & \PseudoGenesInHISTOne{} \\
        HIST2   & \HTwoACodingInHISTTwo{}     & \HTwoAPseudoInHISTTwo{}
                & \HTwoBCodingInHISTTwo{}     & \HTwoBPseudoInHISTTwo{}
                & \HThreeCodingInHISTTwo{}    & \HThreePseudoInHISTTwo{}
                & \HFourCodingInHISTTwo{}     & \HFourPseudoInHISTTwo{}
                & \CodingGenesInHISTTwo{}     & \PseudoGenesInHISTTwo{} \\
        HIST3   & \HTwoACodingInHISTThree{}   & \HTwoAPseudoInHISTThree{}
                & \HTwoBCodingInHISTThree{}   & \HTwoBPseudoInHISTThree{}
                & \HThreeCodingInHISTThree{}  & \HThreePseudoInHISTThree{}
                & \HFourCodingInHISTThree{}   & \HFourPseudoInHISTThree{}
                & \CodingGenesInHISTThree{}   & \PseudoGenesInHISTThree{} \\
        HIST4   & \HTwoACodingInHISTFour{}    & \HTwoAPseudoInHISTFour{}
                & \HTwoBCodingInHISTFour{}    & \HTwoBPseudoInHISTFour{}
                & \HThreeCodingInHISTFour{}   & \HThreePseudoInHISTFour{}
                & \HFourCodingInHISTFour{}    & \HFourPseudoInHISTFour{}
                & \CodingGenesInHISTFour{}  & \PseudoGenesInHISTFour{} \\
        \addlinespace
        Total   & \HTwoACodingGenes{}       & \HTwoAPseudoGenes{}
                & \HTwoBCodingGenes{}       & \HTwoBPseudoGenes{}
                & \HThreeCodingGenes{}      & \HThreePseudoGenes{}
                & \HFourCodingGenes{}       & \HFourPseudoGenes{}
                & \TotalCodingGenes{}       & \TotalPseudoGenes{} \\
        \bottomrule
      \end{tabular}
    \end{table*}


  \section{Histone transcripts}
  %% There is usually reference to histones being the only protein coding
  %% genes transcribed by Pol II that lack poly(A).  However, only Pol II
  %% transcribe protein coding genes so there's no point in making reference
  %% to Pol II.  Other reference dn't even say "only", they may say that
  %% "almost all other protein coding genes" but then have no reference to
  %% what other genes lack poly(A).  I am assuming they are only being
  %% cautious as I can't find information on such genes.


  Canonical core histone transcripts are among the only protein coding messages
  that do not undergo polyadenylation,
  instead carrying a unique stem-loop structure (\tref{tab:typical-histone-differences}).
  They do have a 7--methyl--guanosine cap at the 5' end \citep{MarzluffNatRevGen2008}.

  Canonical histone gene transcription is regulated
  by cell-cycle dependent phosphorylation of the histone-specific NPAT coactivator
  and interaction with the accessory protein FLASH,
  resulting in assembly of histone locus bodies
  coordinating factors responsible for canonical histone transcription and processing
  \citep{MarzluffNatRevGen2008,RattrayMueller2012,Hoefig2014}.
  Variability is observed in canonical histone isoform gene transcription,
  both by analysis of non-polyadenylated transcripts \citep{YangGenomeBiol2011}
  and RNA polymerase II promoter occupancy \citep{Ederveen2011}.

  Overall there is estimated to be a 35 fold increase in mammalian
  canonical histone transcripts during S phase,
  principally as a result of a 10 fold increase in mRNA stabilisation
  via stem-loop dependent mechanisms,
  and a 3--5 fold up-regulation in canonical histone gene transcription \citep{HarrisMCB1991}.

  This post-translation regulation of canonical histone transcripts
  is achieved by a stem-loop encoded in the mRNA 3' untranslated region
  that is recognised by specialised spliceosome-related RNA
  processing and stabilisation complexes \citep{stem-loop-structure}.
  The modal location of annotated stem-loops is \StemLoopStart{} to \StemLoopEnd{} bp after the stop codon
  for the set of all canonical histone genes.
  The sequence logo of all aligned stem-loops
  confirms these are highly conserved in sequence and structure (\fref{fig:stem-loop-seqlogo}).

  The RNA stem-loop structure is bound by the Stem-Loop Binding Protein (SLBP)
  which is up-regulated 10--20 fold during S~phase to stabilise
  histone mRNAs \citep{SLBP-regulation}.
  Immediately downstream of the stem-loop a purine-rich Histone Downstream Element (HDE)
  interacts with U7 snRNA to direct efficient 3' end processing.
  Although this is not an annotated feature of histone genes,
  alignment of the canonical HDE \citep{HDE-sequence} to all canonical histone genes
  shows the modal location of the HDE is 
  \HDEsDistanceFromStemLoop{} bp downstream of the stem-loop.
  The sequence logo calculated from all aligned HDEs 
  confirms that this feature is also highly conserved (\fref{fig:HDE-seqlogo}).

  \begin{figure}
    \centering
    \subtop[Stem-loop]{
      \includegraphics{seqlogo_stem_loops.pdf}
      \label{fig:stem-loop-seqlogo}
    }
    \hfill
    \subtop[HDE]{
      \includegraphics{seqlogo_HDEs.pdf}
      \label{fig:HDE-seqlogo}
    }
    \caption{%
      Sequence logos for A) annotated stem-loops, and 
      B) Histone Downstream Elements (HDEs) identified by homology
      for all canonical histone gene 3' untranslated regions (UTRs).
    }
  \end{figure}

  Stabilisation and processing extend the half life of canonical histone \mbox{mRNAs}
  and contribute to increased histone translation efficiency,
  enabling rapid production of canonical histones to package the newly duplicated genome.

  Although the vast bulk of canonical histone transcripts 
  appear to be regulated by this mechanism,
  1-5\% of canonical histone transcripts
  are found to be 3' polyadenylated \citep{YangGenomeBiol2011}
  and some genes have annotations indicating two transcripts
  differing in whether they have stem-loop or polyadenylation signals (\tref{tab:curation-anomalies}).
  Histone variant transcripts lack a 3' stem-loop and are
  polyadenylated in the same way as almost all other protein coding genes.
  The exception is H2AX, which has alternatively processed transcripts
  exhibiting the stem-loop characteristic of canonical histones
  or a poly(A) tail found on other variant histones \citep{HTwoAX-transcripts,our-H2AX-review}.

  \section{The proteins}
  While some of the histone genes encode the same protein, this is not always the
  case. On the case of H3 and H4 the number of unique proteins are low, \TotalHThree{}
  and \TotalHFour{} respectively. For H2A and H2B this number is much higher, \TotalHTwoA{}
  and \TotalHTwoB{}.

  The numbering of these proteins amino acids also has a catch. The first methionine
  is cleaved after expression and as such, it is never counter when numbering them.

  Most of the PTM are in the histone tails but a few are within the histone fold (H3 T45 for example)

  \subsection{Histone Fold Domain}
    \todo[inline]{Is histone protein structure too much off topic?}

    All histone proteins have a common motif named the Histone Fold Domain (HFD).
    It is composed of a long $\alpha$ helix and two loops on its side that fold
    back two other shorter $\alpha$ helices. This structure allows the histones
    to dimerize with their pair (H3--H4 and H2A--H2B) and form the histone
    octamer. The central $\alpha$ helix of each dimer pair fits in the gap between
    the $\alpha$1 and $\alpha$2 helices crossing the pair $\alpha$2 helix. In
    addition to these, H3 and H2B display a short $\alpha$N and $\alpha$C helix
    on their terminals.

    The HFD is common to all histone and shared with other
    proteins such as the CENP--TWSC and NF--Y complexes. All of these are nuclear
    proteins that interact with DNA. In the case of NF--Y, these even replace
    the H2A--H2b histones during active conditions.

    Whether these proteins actually ever happen to duplicate the nucleosome functions...


  \subsection{H2A}
    \begin{TableAndFigure*}
      \captionof{table}{histone H2A protein consensus}
      \label{tab:H2A-consensus}
      \input{results/table-H2A-align}

      \includegraphics[width=\textwidth]{seqlogo_H2A.pdf}
      \captionof{figure}{WebLogo of all H2A sequences. The bottom line are the amino acids
                         of H2AFJ whose sequence is different from any of the H2A proteins}
      \label{fig:H2A-weblogo}
    \end{TableAndFigure*}

    \todo[inline]{comment on H2AFJ}

  \subsection{H2B}
    \begin{TableAndFigure*}
      \captionof{table}{histone H2B protein consensus}
      \label{tab:H2B-consensus}
      \input{results/table-H2B-align}

      \includegraphics[width=\textwidth]{seqlogo_H2B.pdf}
      \captionof{figure}{WebLogo of all H2B sequences.}
      \label{fig:H2B-weblogo}
    \end{TableAndFigure*}

  \subsection{H3}
    \begin{TableAndFigure*}
      \captionof{table}{histone H3 protein consensus}
      \label{tab:H3-consensus}
      \input{results/table-H3-align}

      \includegraphics[width=\textwidth]{seqlogo_H3.pdf}
      \captionof{figure}{WebLogo of all H3 sequences.}
      \label{fig:H3-weblogo}
    \end{TableAndFigure*}

  \subsection{H4}
    %% these should all be the same but they are not. And there's an H4 in HIST4
    %% which is also present in mouse, which means has been conserved  in evolution
    With the exception of HIST1H4G, all H4 genes seem to encode the same protein. Even HIST4H4 which is the
    only gene on that ``cluster''.
    \begin{TableAndFigure*}
      \captionof{table}{histone H4 protein consensus}
      \label{tab:H4-consensus}
      \input{results/table-H4-align}

      \includegraphics[width=\textwidth]{seqlogo_H4.pdf}
      \captionof{figure}{WebLogo of all H4 sequences.}
      \label{fig:H4-weblogo}
    \end{TableAndFigure*}


  \section{Reproducible research}
\label{sec:reproducible}

	Curation of genomes is an ongoing process, 
	with sequence gaps being completed or corrected,  
	and with annotations being added or updated continuously \addref{some generic curation overview paper}.

	A considerable number of builds and updates of the human genome sequence have been made since 2002 
	so it is unsurprising that 
	\todo{number of sequences added} genes added, 
	\todo{number of sequences removed} removed, 
	and \todo{number of sequences changed} sequences changed  \tref{tab:difference-from-Marzluff02}
	since the last major survey of human canonical histone genes \citep{Marzluff02}.

	Resequencing efforts such as the 1000 genomes initiative have likely  most coding regions have stabilised and 

	In addition to maintaining a current reference, 
	it is also necessary for communities of researchers to feed back information 
	to ensure that reference databases are rich and accurate. 
	In the course of our analyses we have submitted a number of proposals 
	for incorrect, inconsistent or missing annotations of coanonical histone genes. 
	For example, \todo{need a brief example}.
	
	One  process of he continuous curation of genomic databases and it was this
  that led us to publish an updated catalogue of the human histones. The data being
  used in this publication is still being updated. Indeed, it will never stop being
  updated. Since the goal of this publication is an up to date analysis of all histone
  genes, it will become outdated or incorrect once published.
  In addition,
  we are also aware how even the same data can give different results if analysed
  in different ways. For example, \fref{fig:H2A-weblogo} will be different based
  on the alignment algorithm being used such as TCoffee\citep{tcoffee2000},
  ClustalW\citep{clustalw2}, or MultAlin\citep{multalin1988}.

  Since the nature of this publication is a moving target, we have created a set of tools
  to continuously update it based on the latest data that can be found online.
  All sequences were obtained and analysed automatically via a collection of perl
  programs, and we made extensive use of other perl modules, mainly BioPerl\citep{bioperl}.
  When possible, we provided patches to our dependencies, avoiding the infamous
  ``not invented here'' syndrome. This not only makes it easier for colleagues to
  reuse it, whose work may have nothing to do with histones, but also to
  increase its availability.

  At the core of our work is a perl program (bp\_genbank\_ref\_extractor) which
  we incorporated into the Bio-EUtilities distribution since version 1.73. This program
  downloads all sequences --- genomic, transcript, and protein --- related to a search
  result in Entrez gene. Several options have been implemented to select specific genome
  builds, or regions of the genome through offsets from the search results. For example,
  it is possible to download the 500bp upstream of all genes in a search result.

  %% genes may not change a lot but this a good framework to still study the promotor

  For the specific case of human histones, we may have reached a point where the amount
  of changes over time is too small to justify ths work. However, this project was
  designed with reuse in mind, and a catalogue for other organisms can be done with a
  one line change. We have already done this for \textit{Mus musculus} and
  \textit{Gallus gallus} which are of importance in our institute.
  We reckon that our code could also be reused for other families of genes, such as our
  previous catalogue of Snf2 family members\citep{andrew-snf2-catalogue}.

%    We reason that the use of mature software development tools, could be used in other quantitative biology.
  Following the principles of reproducible research
  \citep{reproducible-research-bioinformatics, reproducible-research-law}, we made
  not only all of the source code freely accessible online, but also the instructions
  for automatic builds with SCons\citep{SCons2005}
  (see \url{https://github.com/af-lab/histone-catalog}). The use of such build systems is
  standard practice in software development and is becoming more common in the field
  of computer sciences. But we see this publication with the results in human readable
  form as the goal, not its raw results. Therefore, the final target of our build system is
  the PDF for this publication.

  %% How can we talk about the idea at
  %% http://juretriglav.si/how-scientific-figures-should-work-in-2014/
  %% if we can't cite URL? A very cool idea was mentioned there and deserves
  %% attribution:
  %%
  %%    A great example is a continuous stream of data from an OpenEHR API: why
  %%    wouldn’t that retrospective study on pancreatic cancer survival using
  %%    drug X be a living research object, updating itself as long as the data
  %%    keeps coming?


  Our interest in making this a live and reproducible publication was the basis of
  our choice for the search and download of sequences. Sequences are download from
  RefSeq through Entrez, instead of HAVANA through Ensembl, only because it is much
  easier to make work in other systems. Entrez has made available E-utilities with
  a SOAP interface which is already part of the Bioperl project.
  On the other hand, the Ensembl search is done by a Lucene server which is not publicly
  accessible. The set up of a Lucene instance in a local system is not a straightforward
  solution.
  In addition, the Ensembl perl API is also not part of the Bioperl project, and not
  even available through CPAN. Its dependency on a severely outdated version of BioPerl
  is also a further complication of its installation, to the point that distribution of
  a virtual machine is recommended.
  All this would make reproducibility
  of our data a much more difficult process by other researchers.



  \section{Conclusion}
\label{sec:conclusion}

  Comprehensive curation of the canonical core histone gene family
  reveals significant diversity of isoforms
  and challenges assumptions of bulk chromatin homogeneity.
  It provides a reference for interpreting genomic and proteomic datasets
  and uncovers a number of canonical histones that are targets for future investigation.
  An overview of histone terms and nomenclature
  are included to facilitate clarity in usage of terminology.
  As a dynamic manuscript this catalogue provides a continuously up-to-date overview
  of the human canonical core histone gene family.

  \section{Materials and methods}
\label{sec:matmethods}

  The primary manuscript is generated from a LaTex source
  deriving dynamic data from automatic builds using SCons \todo{SCons version} \citep{SCons2005}.
  Core processing is performed by bp\_genbank\_ref\_extractor
  available in the Bio-EUtilities \todo{BioPerl version} \citep{BioPerl2002}.
  Data is obtained from NCBI RefSeq release \todo{RefSeq release} \citep{PruittRefseq2014}.
  Alignments were generated by T-Coffee \todo{T-Coffee version} \todo{parameters?} \citep{tcoffee2000},
  sequence logos were generated using WebLogo \todo{weblogo version} \todo{parameters?} \citep{weblogo},
  and codons were analysed using SNAP \todo{SNAP version} \citep{KorberSNAP2000}
  and CUSP \todo{EMBOSS version} \citep{Emboss2000}.
  Sequences are represented following Human Genome Variation Society recommended nomenclature \citep{mutnomenclature2003}.
  All scripts, latex sources and current sequence files
  are available in a github repository \url{https://github.com/af-lab/histone-catalog}


  \section{Acknowledgements}
    We are grateful to the Perl community on the \url{irc.freenode.net}
    channels \#perl and \#bioperl,
    especially Altreus, mst, LeoNerd, and pyrimidine.
    Invaluable support was also provided by the \TeX{} stackexchange community,
    in particular David Carlisle.

  \bibliography{references}

  \newpage
  \appendix
  \section{Supplementary Figures}

\beginsupplement

%% Supplementary tables and figures
\onecolumn

\newpage
\captionof{table}{Differences between the current human histone genes and \citep{Marzluff02}.}
\label{tab:difference-from-Marzluff02}
\input{results/table-reference_comparison}

  %% this is a very long table (ctabular) that does not play nice with two columns
  %% so we place it at the end. See why it was a problem at
  %% http://tex.stackexchange.com/questions/101791/float-a-wide-ctabular-in-two-columns
  %% We could split the table in two parts in which case we could use begin{table*}
\newpage
\captionof{table}{Catalogue of all annotated human canonical core histone gene, transcript and protein IDs}
\label{tab:histone-catalogue}
\input{results/table-histone_catalogue}

\begin{figure*}
	\centering
	\subtop[H2A]{
		\includegraphics[width=1.0\textwidth]
		{seqlogo_H2A_cds.pdf}
		\label{fig:H2A-cds-weblogo}
	  }
	\hfill
	\subtop[H2B]{
		\includegraphics[width=1.0\textwidth]
		{seqlogo_H2B_cds.pdf}
		\label{fig:H2B-cds-weblogo}
	}
	\subtop[H3]{
	\includegraphics[width=1.0\textwidth]
		{seqlogo_H3_cds.pdf}
		\label{fig:H3-cds-weblogo}
	}
	\hfill
	\subtop[H4]{
		\includegraphics[width=1.0\textwidth]
		{seqlogo_H4_cds.pdf}
		\label{fig:H4-cds-weblogo}
	}
	\captionof{figure}{Sequence logos for canonical coding histone genes from \tref{tab:histone-catalogue}. 
	showing coding regions of A. H2A, B. H2B, C. H3 and D. H4.}
	\label{fig:histone-gene-variation}
	\end{figure*}

\newpage
\begin{table*}
	\caption{
	Sequence differences between canonical histone genes.
	Average counts of synonymous (Sd) and non-synonymous (Sn) nucleotide substitutions 
	between all pairwise comparisons of canonical histone coding regions by type 
	and average ds/dn ratios for the pairwise comparisons, 
	calculated by SNAP.}
	\label{tab:histone-gene-differences}
	\centering
	\begin{tabular}{r r r r}
		\toprule
		\null & Sd & Sn & ds/dn \\
		\midrule
		H2A	& 52.7 & 5.8 & 75.1 \\
		H2B	& 36.9 & 5.5 & 58.7 \\
		H3	& 57.1 & 1.8 & 214.4 \\
		H4	& 49.2 & 3.9 & 204.7 \\
		\bottomrule
	\end{tabular}
\end{table*}

\newpage
\begin{table*}
	\caption{
	Codon usage percentages for all canonical histone coding regions 
	calculated by CUSP.}
	\label{tab:histone-gene-codonusage}
	\centering
	\begin{tabular}{r r r r}
		\toprule
		\null &  & ds/dn \\
		\midrule
		Ala	&  \\
		Arg	& 36.9 & 5.5 & 58.7 \\
		H3	& 57.1 & 1.8 & 214.4 \\
		H4	& 49.2 & 3.9 & 204.7 \\
		\bottomrule
	\end{tabular}
\end{table*}

\newpage
\captionof{table}{Catalogue of human variant core histone gene, transcript and protein IDs}
\label{tab:variant-catalogue}
\input{results/table-variant_catalogue}

\newpage
\input{results/histone_insanities.tex}


\end{document}
